%=======================================================================
%  Paper V — Section 2 (Data and methods)   FINAL
%
%  Required in the master preamble (once only):
%
%  \usepackage{amsmath,amssymb,booktabs,graphicx,natbib,siunitx}
%  \sisetup{per-mode=symbol, range-phrase={--}, range-units=single,
%           separate-uncertainty=true}
%
%  \newcommand{\potvarn}{\tilde{\sigma}_{U}}
%  \newcommand{\Upot}{U}
%  \newcommand{\Ubar}{\bar{U}}
%  \newcommand{\Ufull}{U^{\mathrm{grav}}_{\mathrm{full}}}
%  \newcommand{\UR}{U^{\mathrm{grav}}_{R}}
%  \newcommand{\Urot}{U^{\mathrm{rot}}}
%  \newcommand{\denshead}{\varrho_{\mathrm{head}}}
%  \newcommand{\densrest}{\varrho_{\mathrm{rest}}}
%  \newcommand{\densR}{\varrho_{R}}
%  \newcommand{\densbulk}{\varrho_{\mathrm{bulk}}}
%  \newcommand{\densneck}{\varrho_{\mathrm{neck}}}
%  \newcommand{\densbody}{\varrho_{\mathrm{body}}}
%  \newcommand{\Ddens}{\Delta\varrho}
%  \newcommand{\Dmass}{\Delta M}
%  \newcommand{\VR}{V_{R}}
%  \newcommand{\Vtot}{V_{\mathrm{tot}}}
%  \newcommand{\Mtot}{M_{\mathrm{tot}}}
%  \newcommand{\ampfac}{\mathcal{A}}
%  \newcommand{\improve}{I}
%  \newcommand{\dCOM}{\Delta x_{\mathrm{COM}}}
%  \newcommand{\xsplit}{x_{\mathrm{split}}}
%  \newcommand{\CV}{\mathrm{CV}}
%  \newcommand{\Sstat}{\mathcal{S}}
%
%  Values superseded from earlier drafts of this section:
%      scan interval        [alpha,beta] -> 300-4000, widened to 300-8000
%      "parabola refinement"             -> removed; none is applied
%      curvature -> rho_rest  1.15 %     -> 1.38 %  (multiplier f/(1-f))
%      rho_rest interval     +/- 23      -> +/- 27 kg/m3
%      Itokawa plane in tab:planes 150 m -> 160.8 m (detector)
%      "Four checks"                     -> Five checks
%
%  All values confirmed. Spin-period provenance and the reproducibility
%  environment are now stated in the text; the only placeholder left in this
%  section is the Zenodo DOI, which set_doi.py fills.
%=======================================================================

\section{Data and methods}
\label{sec:methods}

\subsection{The density model, the mass constraint and the geopotential}
\label{sec:methods:model}

We divide a body into a region $R$ and its complement, assign a uniform
density to each, and hold the total mass at its measured value, as in
Equation~\eqref{eq:massbalance}. One free parameter remains. Writing
$\Ufull$ for the gravitational potential of the whole body at unit
density and $\UR$ for that of the region alone, linearity of the
Newtonian potential in density gives
%
\begin{equation}
U^{\mathrm{grav}}(\mathbf{r})
  = \densrest\,\Ufull(\mathbf{r})
  + \left(\densR - \densrest\right)\UR(\mathbf{r}) ,
\label{eq:superposition}
\end{equation}
%
so that a scan over $\densR$ requires only two evaluations of the
polyhedral gravity field, not one per trial density.

The quantity that relaxation acts on is not $U^{\mathrm{grav}}$ but the
geopotential in the co-rotating frame, which adds a centrifugal term,
%
\begin{equation}
\Upot(\mathbf{r})
  = \densrest\,\Ufull(\mathbf{r})
  + \left(\densR - \densrest\right)\UR(\mathbf{r})
  + \Urot(\mathbf{r}) ,
\qquad
\Urot(\mathbf{r}) = -\tfrac{1}{2}\,\omega^{2}\,d_{\perp}^{2}(\mathbf{r}) ,
\label{eq:geopotential}
\end{equation}
%
with $\omega = 2\pi/P$ and $d_{\perp}$ the perpendicular distance from
the spin axis. We adopt the convention in which the gravitational
potential is positive, so that $\Upot > 0$ over the surfaces considered
here and the normalization by $\lvert\Ubar\rvert$ in
Equation~\eqref{eq:dispersion} is unambiguous.

Equation~\eqref{eq:geopotential} carries the point on which the whole
construction turns, and we state it explicitly because it is easy to
lose. The first two terms scale linearly with density; the third does
not. Were $\Urot$ absent, multiplying every density by a constant would
multiply $\Upot$ by that constant and leave $\potvarn$, a normalized
quantity, unchanged---the objective would be flat in the overall density
scale and the minimization would be empty. It is the
density-independence of $\Urot$ that breaks the scaling and gives the
objective an interior minimum. The same argument applied at fixed total
mass is what gives the two-region problem a determinate solution.

Spin axes are taken from the published pole solutions cited in
Table~\ref{tab:shapes}, and the shape models are used in the body-fixed frame
in which they are supplied, with $z$ along the observed pole. That frame is not
exactly the principal frame even for a uniform body: on Itokawa the axis of
maximum moment of the uniform shape lies $\ang{0.77}$ from $z$. We note a consequence of
the two-region model that is not present in the uniform case: a non-zero
$\Ddens$ displaces the centre of mass, and in general rotates the
principal axes, so that the spin axis of the trial body is not exactly
the $z$-axis of the supplied mesh. On Itokawa the displacement is
$\SI{16.86}{\metre}$ along $x$ at the plane returned by the detector, and
ranges from $14.89$ to $\SI{20.52}{\metre}$ across the sweep of
Table~\ref{tab:itokawasweep}; the axis of maximum moment lies $\ang{0.77}$ from the mesh $z$-axis for the
uniform body and $\ang{1.87}$ for the recovered interior, an induced rotation
of $\ang{1.10}$. The recovered centre of mass computed from the same inertia
calculation lies at $x = \SI{17.22}{\metre}$, which is the $\SI{16.86}{\metre}$
displacement plus the $\SI{0.36}{\metre}$ offset of the centre of figure, so
the two calculations agree.
We hold the spin axis fixed at the shape-model $z$-axis, through the
origin, throughout. To measure what that choice costs we repeated the Itokawa
reference inversion with the rotational potential taken instead about the
recovered axis of maximum moment, through the recovered centre of mass, which
changes both the direction and the position of the axis. The recovered head
density moved from $2800$ to $\SI{2875}{\kilo\gram\per\cubic\metre}$, three
scan steps; the contrast from $1.507$ to $1.560$, an increase of $3.5\%$; and
the improvement from $16.9\%$ to $19.2\%$. The shift in $\denshead$ lies
inside the $\pm\SI{129}{\kilo\gram\per\cubic\metre}$ curvature interval of
Section~\ref{sec:itokawa}, and the change in contrast inside the $6.0\%$
error budget given there, so the choice does not alter the result; but it is
not negligible, and it enters that budget as a further term of comparable
size. We do not adopt the recomputed axis as the reference, because the
observed pole is what a real rotation defines and an interior whose principal
axis departs from it is, by that measure, less rather than more consistent
with the data.

Potentials are computed with the Werner--Scheeres formulation
\citep{werner1996,tsoulis2012} as implemented in
\texttt{polyhedral-gravity} version~\texttt{3.3.1}
\citep{gravitylib},
at field points placed at each facet centroid and displaced inward along
the outward normal by $10^{-4}$ of the bounding-box diagonal. The
displacement avoids the removable singularity of the Werner--Scheeres
expressions at a field point lying on a facet, but it is not physically
neutral: the same fractional offset is $\SI{6.5}{\centi\metre}$ on
Itokawa and $\SI{3.8}{\metre}$ on Eros, and it shifts the potential by
$\delta U \simeq g\,d$, of order $0.04\%$ of $\lvert\Ubar\rvert$ on both.
That is comparable to the dispersion improvements reported for the
weak-signal bodies in Section~\ref{sec:nulls}, so we repeated the Itokawa
reference run and the Eros run at Himeros---the strong case and the
best-resolved weak one---with the offset scaled by $10$ and by $0.1$. The
recovered contrasts did not move at the resolution of the density scan. The
improvements moved by at most $0.29$ percentage points on Itokawa and
$0.007$ on Eros, in both cases about $1.7\%$ of the improvement itself, and
in the same direction, a larger offset giving a larger improvement.

The normalized geopotential dispersion follows Paper~IV:
%
\begin{equation}
\potvarn = \frac{\sqrt{\sum_i A_i (\Upot_i - \Ubar)^2 \big/ \sum_i A_i}}
                {\lvert \Ubar \rvert},
\qquad
\Ubar = \frac{\sum_i A_i \Upot_i}{\sum_i A_i},
\label{eq:dispersion}
\end{equation}
%
with $A_i$ the facet areas. The sums run over the facets of the
\emph{whole body}, not of the region: the region enters only through
$\UR$ in Equation~\eqref{eq:geopotential}. Where facets are excluded from
the surface---on Arrokoth, those buried inside the opposite shell
(Section~\ref{sec:methods:shapes})---both $\Ubar$ and $\sum_i A_i$ are
recomputed over the retained set. Equation~\eqref{eq:dispersion} is the
same quantity that \citet{richardson2014} plot as the normalized
potential standard deviation. We report the improvement over the uniform
case,
%
\begin{equation}
\improve = \frac{\potvarn(\text{uniform}) - \potvarn(\densR^{\ast})}
                {\potvarn(\text{uniform})},
\label{eq:improvement}
\end{equation}
%
and the two derived quantities that will matter in
Section~\ref{sec:constrains}: the excess mass and the resulting
displacement of the centre of mass from the centre of figure,
%
\begin{equation}
\Dmass = (\densR - \densrest)\,\VR ,
\qquad
\dCOM = \frac{\Dmass\,(x_R - x_{\mathrm{cf}})}{\Mtot},
\label{eq:excessmass}
\end{equation}
%
where $x_R$ and $x_{\mathrm{cf}}$ are the $x$-coordinates of the
centroids of the region and of the whole body.
Equation~\eqref{eq:excessmass} is exact, not a small-anomaly
approximation: it follows from the definition of the centre of mass of a
two-density body with the complement's centroid eliminated through
Equation~\eqref{eq:massbalance}.

\subsection{Two consequences of the mass constraint}
\label{sec:methods:identity}

Two consequences of Equation~\eqref{eq:massbalance} should be stated
before any result rests on them. The first is the
identity~\eqref{eq:excessidentity}, an exact affine relation between
$\Dmass$ and $\densrest$ at fixed total mass and volume. The two are not
independent quantities: a statement that one is invariant under a change
of dividing plane is the same statement about the other. Differentiating,
%
\begin{equation}
\frac{\delta \Dmass}{\Dmass}
= -\,\ampfac\,\frac{\delta \densrest}{\densrest} ,
\label{eq:excessamplify}
\end{equation}
%
with $\ampfac$ as defined in Equation~\eqref{eq:amplification}. For the
Itokawa reference solution
$\densbulk = \SI{2012.87}{\kilo\gram\per\cubic\metre}$ and
$\densrest = \SI{1858.1}{\kilo\gram\per\cubic\metre}$ give
$\ampfac = 12.0$; across the thirteen planes of
Table~\ref{tab:itokawasweep} it ranges from $11.34$ to $12.70$ with a
mean of $11.77$, and over the nine admissible planes its mean is $11.90$.
Fractional ranges in $\Dmass$ and in $\densR$ are therefore not directly
comparable, and Section~\ref{sec:constrains} reports $\densrest$
alongside both. The identity can be seen operating in the output as well
as in the derivation: the ratio of the measured coefficients of
variation, $\CV(\Dmass)/\CV(\densrest) = 2.946/0.2506 = 11.76$ over all
thirteen planes and $3.00/0.252 = 11.90$ over the admissible nine,
reproduces the corresponding mean $\ampfac$ in both cases.

The second consequence is that $\ampfac$ is itself a poorly conditioned
quantity, because $\densbulk - \densrest$ is a difference of two
comparable numbers. Differentiating
Equation~\eqref{eq:amplification} at fixed $\densbulk$,
%
\begin{equation}
\frac{\delta \ampfac}{\ampfac}
 = \frac{\densbulk}{\densbulk - \densrest}\,
   \frac{\delta \densrest}{\densrest}
 \;\simeq\; 13\,\frac{\delta \densrest}{\densrest}
\quad\text{on Itokawa,}
\label{eq:ampsensitivity}
\end{equation}
%
so that a shift of $0.13\%$ in the background density moves $\ampfac$ by
$1.7\%$. This is why $\ampfac$ varies by $\pm5.6\%$ across a sweep over
which $\densrest$ varies by $\pm0.434\%$, and why any quoted value must
carry the plane at which it was evaluated.

Two region shapes are used. The \emph{head} model takes $R$ to be the
material beyond a plane, $x > \xsplit$. The \emph{neck} model takes $R$
to be a slab, $\lvert x - x_{\mathrm{neck}}\rvert < w/2$. The two are
alternative hypotheses about where the anomaly sits, not independent
constraints, and they overlap whenever $\xsplit$ falls inside the slab.
Section~\ref{sec:itokawa:ambiguity} shows that on Itokawa the two return
background densities differing by more than three times the combined
uncertainty on either, so the choice between them is a modelling decision
with consequences, not a matter of presentation.

\subsection{Region construction by exact plane clipping}
\label{sec:methods:clip}

Each region is built as a closed, consistently oriented triangle mesh.
Every facet crossing the plane is clipped against the half-space by the
reentrant-polygon algorithm of \citet{sutherland1974}; since the subject
polygon is a triangle and the clip region is a single half-space, the
clipped facet is a triangle or a quadrilateral and the general reentrant
machinery reduces to one pass. Intersection points are snapped exactly
onto the plane, and the opening is closed by fanning each boundary edge,
traversed in reverse, to an apex lying in the plane. Fan triangulation of
a planar loop gives the correct signed area from any coplanar apex, so
non-convex cross-sections need no special treatment. The slab of the neck
model is obtained by clipping twice.

Two details of the implementation are load-bearing. An intersection point
is created once per \emph{edge}, not once per facet, so that the two
facets sharing a cut edge reference the same vertex index; created per
facet, they would reference two distinct indices at the same position and
the clipped surface would be torn along the cut---a mesh that still
encloses the correct volume, since that is a sum of signed tetrahedra,
but is not a closed surface. And the cut boundary is traced into separate
loops, each capped to its own apex. A cross-section can be multiply
connected, or the region can fall into more than one piece, in which case
fanning every boundary edge to a single apex would stitch unrelated loops
together. Both conditions arise in this study: the Arrokoth shape model
is supplied as two unjoined shells, so a plane beyond its neck leaves the
region in two pieces (Section~\ref{sec:methods:shapes}).

\subsection{Why the obvious construction fails}
\label{sec:methods:cone}

An earlier version of this work closed each region differently, by
fanning every boundary edge to a single interior point $O$. That mesh is
watertight, and its volume agrees with the sum of the underlying
tetrahedra to machine precision. It nevertheless encloses the wrong
solid: the radial cone over the retained surface patch, with apex at
$O$---for a sphere, a spherical sector---rather than the plane-cut lobe
intended.

\begin{figure}[htbp]
\centering
\includegraphics[width=0.86\textwidth]{fig01_cone_vs_lobe.pdf}
\caption{The two solids, on a unit sphere cut at $x=0.5$. Top: the
plane-cut lobe intended by the density model, of volume $0.654498$ with
centroid at $0.675000$. Bottom: the radial cone---here a spherical
sector---produced by closing the region to an interior point $O$, of volume
$1.047198$ with centroid at $0.562500$. Both are watertight and both have
exact volumes; the sector has $1.600$ times the volume and its centroid lies
$0.1125R$ further in. Of the five checks of
Section~\ref{sec:methods:verify} only the vertex predicate distinguishes
them.}
\label{fig:cone}
\end{figure}

A unit sphere cut at $x=0.5$ makes the difference exact
(Figure~\ref{fig:cone}). The intended spherical cap has volume
$\pi h^{2}(3R-h)/3 = 0.654498$ and centroid
$3(2R-h)^{2}/[4(3R-h)] = 0.675000$; the spherical sector has volume
$\tfrac{1}{3}R\,A_{\rm cap} = 2\pi R^{2}h/3 = 1.047198$ and centroid
$\tfrac{3}{8}(2R-h) = 0.562500$. The constructions differ by a factor
$1.600$ in volume and by $0.1125\,R$ in centroid. On Itokawa the
consequence was not a small bias but a reversal: raising the density of a
sector adds mass along a wedge reaching back to the centre of the body
rather than in the head, and the recovered contrast came out inverted.

\subsection{Degenerate faces and numerical tolerance}
\label{sec:methods:tol}

Vertices lying within a tolerance of the cutting plane are classified as
\emph{on} it, snapped exactly onto it, and never duplicated as
intersection points; only strict sign changes generate an intersection.
Without this, a vertex sitting on the plane makes both of its edges
register as crossings, producing two coincident points and hence
zero-area triangles on both the surface and the cap. A single null-normal
facet makes the polyhedral evaluation return NaN at \emph{every} field
point, not merely near the plane, so the failure is total rather than
local. The tolerance is $10^{-10}$ of the bounding-box diagonal.

At that tolerance a vertex a few millimetres from the plane on a
kilometre-scale body still produces a sliver, which the gravity library
flags as numerically unstable. We tested whether such slivers matter by
repeating runs with the tolerance raised by four orders of magnitude,
which removes them. On 67P the minimum facet area rose from
$2.381\times10^{-12}$ to $3.660\times10^{-8}$~km$^{2}$ while the
recovered contrast, the improvement and the excess mass changed in the
seventh significant figure or beyond; on Itokawa, the body from which the
principal result of this paper comes, the recovered $\denshead$ was
unchanged at every plane of Table~\ref{tab:itokawasweep} and $\densrest$
was unchanged to four decimal places;
on a synthetic body every printed digit was unchanged. The slivers are
harmless, as the physics requires---their contribution to
Equation~\eqref{eq:dispersion} scales with their area---but this was
established by measurement rather than assumed.

\subsection{Verification}
\label{sec:methods:verify}

Five checks guard the calculation. Three of them can fail in ways the
others cannot, and that is the reason for having five.

\emph{Analytic.} The reference solid is an icosphere of $20{,}480$ facets
inscribed in the unit sphere.
The clipped cap has volume $0.653931$ against the closed form $0.654498$
($-0.087\%$) and centroid $0.674936$ against $0.675000$ ($-0.009\%$); a
slab between $x=\pm0.2$ has volume $1.239437$ against $1.239882$
($-0.036\%$). These residuals are the polyhedral approximation of the
sphere itself---the whole-body volume of the same mesh falls short of
$4\pi/3$ by $0.05\%$, and an inscribed polyhedron necessarily
underestimates---not error in the clipping, which is exact for the given
polyhedron. That the cap and slab deficits bracket the whole-body
deficit, and that all three carry the same sign, is what identifies them
as discretization rather than clipping error.

\emph{Complement.} $V(x>c) + V(x<c)$ must equal the body volume. The
residual on Itokawa is $2.0\times10^{-14}$ per cent, that is
$2.0\times10^{-16}$ in relative terms. That is below the double-precision
epsilon of $2.2\times10^{-16}$, and below the $\sqrt{N}\,\epsilon \approx
7\times10^{-14}$ that an uncorrelated sum of $10^{5}$ signed tetrahedra would
give. The explanation is that the two half-space meshes share most of their
facets with the whole-body mesh, so the round-off is common to both sides of
the comparison and cancels rather than accumulating. The test is therefore
weaker than its magnitude suggests: it confirms that the clipping partitions
the tetrahedral sum, which is why Section~\ref{sec:methods:cone} adds the
predicate test.

\emph{Predicate.} Every vertex of the region must satisfy its defining
inequality. The plane-cut lobe passes with zero violation. The spherical
sector of Section~\ref{sec:methods:cone} fails outright, its apex lying
at the centre of the body, with a violation of $0.5$ in units where the
plane sits at $x=0.5$.

\emph{Manifold.} Every directed edge must appear exactly once and its
reverse must be present, and the Euler characteristic
$\chi = n_V - n_E + n_F$ must equal $2$ for each connected component,
where $n_V$, $n_E$ and $n_F$ count vertices, edges and facets. This is
independent of the other checks, and necessarily so: the
divergence-theorem volume is a sum of signed tetrahedra and requires only
that the oriented facets tile the boundary with correct signs, not that
they are stitched together. A surface torn along the cut therefore passes
the analytic, complement and predicate tests while not being a closed
surface at all. The criterion is stated per component rather than as a
single $\chi = 2$, because a region cut from a body supplied as several
unjoined shells is legitimately in several pieces.

\emph{Mass balance.} Every solution must satisfy
$\Dmass = (\densbulk - \densrest)\Vtot$ and, independently,
$\Dmass = \Ddens\,\VR$. Both hold to the printed precision on all
thirteen rows of Table~\ref{tab:itokawasweep}. The check is enforced by
construction and cannot fail unless the region volume passed to the mass
constraint differs from the one passed to the gravity evaluation---which
is exactly the failure mode of Section~\ref{sec:methods:cone}, where the
two differed by a factor $1.600$. It is therefore the layer that would
have caught that error at the level of the output rather than of the
geometry, and it is the reason we tabulate $\densrest$ and $\Dmass$ side
by side rather than quoting either alone.

The first three checks run on every region and abort the model if they
fail. The manifold check reports rather than aborts, and distinguishes a
defect inherited from the source mesh from one introduced by the
clipping; the distinction matters because a volume and a potential are
unaffected by tearing, so a defect that the source already carries cannot
be repaired here and should not stop the calculation. The mass-balance
audit runs on every solution and is reported with the results.
Section~\ref{sec:methods:shapes} records which shape models carry
topological defects.

\subsection{External validation of the geometry}
\label{sec:methods:external}

The preceding subsections test the construction against solids of our own
choosing. A stronger check is available for 67P, whose small lobe has
been measured independently. Our clipped region holds $26.6\%$ of the
body volume against the $27\%$ of \citet{jorda2016}, at the plane the
detector of Section~\ref{sec:methods:neck} returns with nothing tuned.
The 67P model is a single connected surface, so producing that figure
requires the clipping to do real geometric work.

Three qualifications belong with that number, and we state them rather
than let the agreement stand unexamined. First, \citet{jorda2016}
decompose the nucleus into two lobes \emph{and} a neck, the neck
accounting for about $7\%$ of the total volume and assigned to neither
lobe, whereas a single plane necessarily assigns the entire neck to one
side. Our figure therefore agrees with theirs only because the detected
plane falls at the far edge of the neck; displacing it across the neck
moves the fraction by $10.7$ percentage points, from $31.80\%$ at
$\xsplit = \SI{0.70}{\kilo\metre}$ to $21.14\%$ at $\SI{1.20}{\kilo\metre}$,
a gradient of $-21.2$ points per kilometre. Sweeping the $7\%$ that
\citet{jorda2016} assign to the neck therefore corresponds to about
$\SI{330}{\metre}$ of plane travel. We report that sensitivity in
Table~\ref{tab:planes} so that the reader may judge how much of the
agreement is geometric fidelity and how much is plane placement.
Second, the comparison is between a volume fraction and a volume
fraction; the mass fractions quoted by \citet{jorda2016} coincide with
them only under the assumption of uniform density, which is the
assumption this paper exists to test. Third, our $\Vtot$ of
$\SI{18.5912}{\cubic\kilo\metre}$ comes from an SPG model, whereas the
lobe decomposition of \citet{jorda2016} is based on the SPC SHAP5 model
with $\Vtot = \SI{18.8+-0.3}{\cubic\kilo\metre}$; the $1.1\%$ difference
in total volume is smaller than the effect just described but is not
zero.

The corresponding comparison for Arrokoth is weaker, and we report it as
such. Our region holds $66.33\%$ against $66.30\%$ from
\citet{porter2024}, but the model is supplied as two unjoined shells and
the larger of them is $66.30\%$ of the summed component volumes by itself
(Section~\ref{sec:methods:shapes}). The agreement is therefore largely a
property of the model as supplied rather than an independent test of the
clipping, and only the 67P comparison should be read as validation.

\subsection{Locating the dividing plane}
\label{sec:methods:neck}

The cross-sectional half-width is binned along $x$ and a saddle sought
within it. A neck is a \emph{local} minimum flanked by a hump on either
side, not merely the smallest half-width in the interior of the profile.
The distinction is not academic: on a body that tapers at both ends the
smallest interior half-width lies on a flank. On Eros the
smallest-interior-minimum rule would select
$x = \SI{-10.76}{\kilo\metre}$, on a shoulder, whereas the saddle rule
adopted here selects Himeros at $x = \SI{+4.47}{\kilo\metre}$.
% Confirmed against the run records: -10.7619 km is returned by the rejected
% smallest-interior-minimum rule, +4.4698 km by the saddle rule used here.
Section~\ref{sec:nulls} shows what the substitution costs.

Candidates are therefore detected on runs of equal profile value---a
strict point test finds nothing on a symmetric profile, where two samples
straddling the minimum carry equal values---and each is scored by
%
\begin{equation}
p = \frac{\min(w_{\mathrm{left}}, w_{\mathrm{right}}) - w_{\mathrm{neck}}}
         {\max_x w(x)},
\label{eq:prominence}
\end{equation}
%
with $w_{\mathrm{left}}$ and $w_{\mathrm{right}}$ the flanking maxima.
The most prominent candidate is selected provided $p \ge 0.02$;
otherwise the routine returns no neck and the plane must be supplied
explicitly. Where the selected saddle differs from the absolute interior
minimum, the rejected position is reported.

Measured prominences separate the bodies cleanly: Arrokoth $0.490$, Eros
$0.216$, Itokawa $0.170$, 67P $0.091$, and Bennu $0.0004$, which is
refused. The gap between the lowest accepted value and the highest
rejected one spans more than two orders of magnitude, so the threshold is
not a fine judgement. Figure~\ref{fig:profiles} shows the three cases the
detector must distinguish.

\begin{figure}[htbp]
\centering
\includegraphics[width=\textwidth]{fig02_profiles.pdf}
\caption{Cross-sectional half-width profiles for the three cases the saddle
detector must distinguish. Itokawa (top): a clear saddle that is also the
smallest interior half-width, accepted at a prominence of $0.170$, one of
three local minima. Eros (middle): the true saddle at Himeros, prominence
$0.216$, while the smallest interior half-width lies on a flank at
$x=-10.76$~km and is rejected---taking it instead costs what
Section~\ref{sec:nulls} reports. Bennu (bottom): no saddle at all, the
largest of three local minima having a prominence of $0.0004$, and the
detector declines to return a plane.}
\label{fig:profiles}
\end{figure}

Bennu is retained in the study precisely because
the detector refuses it: it supplies a case in which the two-region model
is applied to a body with no neck at all, and Section~\ref{sec:nulls}
reports the inversion at a plane through the centre of figure, chosen by
us rather than by the detector, to show what the objective returns when
the model has no geometric warrant.
% Confirmed: the Bennu plane is x = 0 exactly, imposed by us, enclosing
% 50.24 per cent of the volume; the detector returns no plane on this body.

Table~\ref{tab:planes} collects the plane used for every body, together
with its provenance, so that no reported result depends on a plane whose
origin is stated only in passing. The two Itokawa entries are listed
separately and should not be conflated: the detector returns
$\xsplit = \SI{160.8}{\metre}$, enclosing $16.43\%$ of the volume, while
the plane of \citet{kanamaru2019} lies at $\SI{150}{\metre}$ and encloses
$17.89\%$. The first is the reference solution of this paper; the second
is used only where a direct comparison with published values is intended,
and every quoted number states which of the two it belongs to.

\begin{table}[htbp]
\centering
\small
\caption{Dividing planes used in this work. ``Detector'' means the saddle
returned by Equation~\eqref{eq:prominence} with no tuning; ``imposed''
means supplied by us, with the reason given. $f = \VR/\Vtot$ is the
volume fraction of the region. The two Itokawa rows are distinct planes
separated by $\SI{10.8}{\metre}$ and are kept apart throughout the paper.}
\label{tab:planes}
\setlength{\tabcolsep}{4pt}
\begin{tabular}{llrrrl}
\toprule
Body & Model & $\xsplit$ & $p$ & $f$ (\%) & Provenance \\
\midrule
Itokawa  & head & \SI{160.8}{\metre}      & $0.170$  & $16.43$ & detector; reference solution \\
Itokawa  & head & \SI{150}{\metre}        & ---      & $17.89$ & imposed; \citet{kanamaru2019}, comparison only \\
Itokawa  & neck & \SI{118.7}{}--\SI{202.9}{\metre} & ---   & $12.05$ & imposed; slab, Section~\ref{sec:itokawa:ambiguity} \\
Eros     & head & \SI{+4.47}{\kilo\metre} & $0.216$  & $34.6$  & detector (Himeros) \\
Bennu    & head & \SI{0}{\kilo\metre}    & $0.0004$ & $50.2$  & imposed; detector refuses \\
67P      & head & \SI{+0.952}{\kilo\metre} & $0.091$  & $26.6^{a}$  & detector \\
Arrokoth & head & \SI{-4.786}{\kilo\metre} & $0.490$  & $66.33$ & detector \\
\bottomrule
\end{tabular}
\end{table}
\par\smallskip{\footnotesize
\textsuperscript{a}\,Across the 67P neck $\mathrm{d}f/\mathrm{d}\xsplit =
-21.2$ percentage points per kilometre; $f$ runs from $31.80\%$ at
$\SI{0.70}{\kilo\metre}$ to $21.14\%$ at $\SI{1.20}{\kilo\metre}$. The
Itokawa neck slab has width $w = \SI{84.2}{\metre}$.}

\subsection{The minimization and its sensitivity interval}
\label{sec:methods:minimisation}

The single free parameter of Equation~\eqref{eq:massbalance} is scanned
rather than solved for, because the objective is cheap once $\Ufull$ and
$\UR$ are in hand: by Equation~\eqref{eq:geopotential} each trial density
costs one vector addition. We scan $\densR$ on a uniform grid of step
$\SI{25}{\kilo\gram\per\cubic\metre}$ over an interval set for each body
relative to its own bulk density, since a range appropriate to a rocky
asteroid would miss the minimum entirely on a cometary nucleus or a
Kuiper Belt object, and report the grid point at which $\potvarn$ is
least. On Itokawa the interval is $300$--$4000$~kg\,m$^{-3}$, widened to
$300$--$8000$~kg\,m$^{-3}$ for the outer planes of the sweep of
Section~\ref{sec:itokawa:sweep}; because the two intervals share the same
grid, widening adds candidate points without displacing existing ones.
The intervals used elsewhere are
$500$--$5000$ on Eros, $200$--$3000$ on Bennu, $100$--$4000$ on both 67P
models, and $25$--$3000$~kg\,m$^{-3}$ on Arrokoth at each of the three assumed
bulk densities, with grid steps of $25$, $10$, $5$ and $5$~kg\,m$^{-3}$
respectively.

A minimum returned at either end of the scan interval is flagged as a
boundary minimum and the interval widened until an interior minimum is
obtained; no boundary minimum is reported as a result. Every solution in
this paper satisfies that condition, and Section~\ref{sec:failures}
describes the one case in which the diagnostic fired during development.

The upper limit of the scan also bounds the range of dividing planes the
sweep can probe. Holding $\Dmass$ at its Itokawa sweep value of
$2.80\times10^{9}$~kg, a region small enough to require $\densR$ at the cap
of $\SI{8000}{\kilo\gram\per\cubic\metre}$ has a volume fraction of
$\Dmass/[(8000-\densrest)\Vtot] = 2.56\%$, which the sweep reaches at
$\xsplit \approx \SI{0.257}{\kilo\metre}$. No plane in the sweep comes
within $\SI{17}{\metre}$ of that, but the figure matters in
Section~\ref{sec:plane}, where an extrapolation beyond the sweep would land
close to it.

No sub-grid refinement is applied, and the reported densities are
therefore quantized at that step. We state this rather than leave it to
be inferred, because it is visible in the results: in
Table~\ref{tab:itokawasweep} the planes at $\SI{160.0}{\metre}$ and
$\SI{160.8}{\metre}$, separated by $\SI{0.8}{\metre}$, return the same
$\denshead$ but different $\densrest$, and at two places $\dCOM$
decreases as the plane advances even though $x_R - x_{\mathrm{cf}}$
increases monotonically throughout. The size of the effect follows from
the mass constraint. A half-step error of
$\SI{12.5}{\kilo\gram\per\cubic\metre}$ in $\denshead$ propagates as
$[f/(1-f)]\times\SI{12.5}{\kilo\gram\per\cubic\metre}
= \SI{2.4}{\kilo\gram\per\cubic\metre}$, or $0.13\%$, against an observed
half-range in $\densrest$ of $\SI{8.05}{\kilo\gram\per\cubic\metre}$
($0.434\%$); in standard-deviation terms, a uniformly distributed
quantization error contributes
$\SI{25}{\kilo\gram\per\cubic\metre}/\!\sqrt{12}$ in $\denshead$ and
hence $\SI{1.4}{\kilo\gram\per\cubic\metre}$ in $\densrest$, against an
observed $\SI{4.65}{\kilo\gram\per\cubic\metre}$. Roughly one tenth of
the observed variance in $\densrest$ is therefore numerical rather than
physical, and the quoted plane-placement scatter is an upper bound on the
physical effect.
% NOTE: Equation (3) is affine in rho_R once the mass constraint is
% substituted, so sigma_U^2 is a ratio of quadratics in rho_R and the
% stationarity condition is linear, with a single root. Solving it in
% closed form would remove the grid step from the budget entirely and
% would also give the curvature interval as the two roots of a quadratic
% rather than as a fitted quantity. This is left for a future revision.

The objective is deterministic: there is no noise model, and repeated
evaluation returns identical values. The interval we quote alongside
$\densR^{\ast}$ is therefore \emph{not} a confidence interval, and we do
not call it one. It is a \emph{curvature sensitivity interval}: the range
of $\densR$ over which $\potvarn$ rises by no more than $0.60\%$ above
its minimum, the threshold being the plateau width between the two finest
meshes as reported in Paper~IV and reproduced in
Appendix~\ref{app:envelope}. Recomputing that plateau for the present study
gives a tighter figure: between the two finest meshes the dispersion moves by
$0.008\%$ on Itokawa and at most $0.098\%$ on any of the four resolution
bodies, with either decimator. We retain $0.60\%$ because it is the published
value and because a wider threshold gives a wider interval, so the sensitivity
intervals quoted throughout are conservative by roughly a factor of six.
It answers the question ``how far can the density move before the
objective would have preferred a different mesh?'' rather than ``what is
the standard error of the estimate?'', and on Itokawa it spans
$\pm\SI{130}{\kilo\gram\per\cubic\metre}$.

Because the mass constraint ties the two densities, that interval
propagates to the background density as
$\delta\densrest = [f/(1-f)]\,\delta\denshead$, which at $f = 0.1643$
gives $\SI{25.6}{\kilo\gram\per\cubic\metre}$, or $1.38\%$. This is three
times the plane-placement half-range of $0.434\%$ and is the dominant
term in the budget; any statement of the form ``the background density is
stable to four parts in a thousand'' is a statement about plane placement
alone and must not be read as an uncertainty. Combining the two in
quadrature at the reference plane,
%
\begin{multline}
\densrest = 1858 \pm \SI{27}{\kilo\gram\per\cubic\metre},
\qquad
\Dmass = (2.75 \pm 0.48)\times10^{9}~\mathrm{kg}
       = 7.7 \pm 1.3\%~\text{of}~\Mtot ,
\qquad \\
\dCOM = 16.9 \pm \SI{2.9}{\metre} .
\label{eq:budget}
\end{multline}
%
The background density differs from the bulk density by
%
\begin{equation}
\Sstat \;\equiv\;
\frac{\lvert \densbulk - \densrest \rvert}{\sigma(\densrest)}
\;=\; 5.8 ,
\label{eq:sstat}
\end{equation}
%
times its own combined uncertainty. That ratio, and not the scatter of
$\densrest$ alone, is what distinguishes a detection from a uniform body:
for a uniform body $\densrest \equiv \densbulk$ holds identically and the
scatter would be zero, so stability alone is evidence of nothing.
$\Sstat$ is used in Section~\ref{sec:itokawa:ambiguity} to compare region
shapes and in Section~\ref{sec:nulls} to compare bodies.

\subsection{Shape models, masses and spin states}
\label{sec:methods:shapes}

\begin{table}[htbp]
\centering
\scriptsize
\setlength{\tabcolsep}{4pt}
\caption{Shape models and parameters. Facets are those used after
decimation, where any was applied. Arrokoth has no measured mass; three
assumed bulk densities were run. The $\chi$ column is the Euler
characteristic of the mesh as used. $\densbulk = \Mtot/\Vtot$ is listed
beside the published value because it enters
Equation~\eqref{eq:amplification} with the amplification of
Equation~\eqref{eq:ampsensitivity}.}
\label{tab:shapes}
\small
\setlength{\tabcolsep}{4.5pt}
\begin{tabular}{lrrrrrrrr}
\toprule
Body & Native & Used & $\chi$ & $\Vtot$ & $\Mtot$ (kg)
     & $\densbulk$ & Published $\densbulk$ & $P$ (h) \\
\midrule
Itokawa  & $49{,}152$      & $49{,}152$ & $2$ & $0.0177856$~km$^3$ & $3.58\times10^{10}$   & $2012.9$ & $1950\pm140$ & $12.1324$ \\
Eros     & $49{,}152$      & $49{,}152$ & $2$ & $2506.39$~km$^3$   & $6.687\times10^{15}$  & $2668$   & $2670\pm30$  & $5.27$ \\
Bennu    & $196{,}608$     & $50{,}000$ & $2$ & $0.061553$~km$^3$  & $7.330\times10^{10}$  & $1191$   & $1190\pm13$  & $4.296061$ \\
67P      & $1{,}309{,}996$ & $50{,}000$ & $2$ & $18.5912$~km$^3$   & $9.982\times10^{12}$  & $537$    & $537.8\pm0.7$ & $12.4043$ \\
Arrokoth & $40{,}960$      & $40{,}960$ & $4$ & $4123.81$~km$^3$   & assumed               & ---      & $155$--$600$ & $15.92$ \\
\bottomrule
\end{tabular}
\end{table}
% Published densities: Itokawa 1950+/-140 (Abe et al. 2006); Eros 2670+/-30
% (Yeomans et al. 2000); Bennu 1190+/-13 (Scheeres et al. 2019); 67P
% 537.8+/-0.7, 3-sigma, from the SHAP7 volume of 18.56 km^3 (Preusker et al.
% 2017), against 532+/-7 from SHAP5 (Jorda et al. 2016); Arrokoth 235, 1-sigma
% range 155-600 (Keane et al. 2022). Arrokoth period 15.92+/-0.02 h (Spencer
% et al. 2020); 67P initial period 12.4041+/-0.0001 h (Jorda et al. 2016).
Spin periods and poles are taken from the following. Itokawa,
$\SI{12.1324(1)}{\hour}$ from the ground-based lightcurve campaign of
\citet{nishihara2019}, which refines the $\SI{12.132(5)}{\hour}$ of
\citet{kaasalainen2003}; Eros, $\SI{5.27025547}{\hour}$ with the pole of
\citet{miller2002}, of which the $\SI{5.27}{\hour}$ used here is a rounding
that changes the rotational term by $8\times10^{-5}$ in relative terms;
Bennu, $\SI{4.296061}{\hour}$, the value determined in flight
\citep{hergenrother2019}, against $\SI{4.297461}{\hour}$ from the
pre-encounter radar and lightcurve solution \citep{nolan2013}---the nucleus
is spinning up, and the two differ by five seconds, which changes the
rotational term by $7\times10^{-4}$; 67P and Arrokoth as given above. The
shape models themselves are identified by their published references rather
than by archive identifiers, which the deposit records
(Section~\ref{sec:conclusions}, data availability): \citet{gaskell2008} for
Itokawa and Eros, the OLA v20 product of \citet{barnouin2019} for Bennu,
SHAP2 and SPC~2017 for 67P \citep{preusker2017,jorda2016}, and
\citet{porter2024} for Arrokoth.
% Confirmed: the Arrokoth V_tot of 4123.81 km^3 is the sum of the two
% components, 1389.53 and 2734.28 km^3; they do not overlap, the intersection
% being 0.044 per cent of the total (Section 2.7).

Masses are taken from \citet{abe2006} for Itokawa, \citet{yeomans2000}
for Eros, \citet{scheeres2019} for Bennu, whose $7.329\pm0.009\times10^{10}$~kg
we round to $7.330$, and \citet{patzold2016} for 67P. The published 67P
density of $537.8\pm0.7$~kg\,m$^{-3}$ is the SHAP7 value of
\citet{preusker2017}; the SHAP5 volume of \citet{jorda2016} gives
$532\pm7$, and our $537$ from a decimated model of $18.5912$~km$^{3}$ lies
between them. The 67P rotation period is the initial value, which held from
arrival until October 2014; \citet{jorda2016} give $12.4041\pm0.0001$~h for
it, $0.7$~s from the value used here. The period then lengthened to
$12.4304$~h by May 2015 and fell to $12.305$~h just before the August 2015
perihelion, so the epoch matters: the change is several hundred times the
quoted uncertainty.
Arrokoth has no mass determination \citep{spencer2020}, and we run three
assumed bulk densities, $250$, $400$ and
$\SI{500}{\kilo\gram\per\cubic\metre}$, with a fourth at $235$, spanning the range discussed by
\citet{mckinnon2020} and \citet{spencer2020}. The published constraints do not
agree with one another. \citet{keane2022} infer $235$~kg\,m$^{-3}$ with a
$1\sigma$ range of $155$--$600$; \citet{spencer2020} require more than
$290$ if the neck is not in tension; and \citet{mckinnon2020} find
$250$--$500$ for a comet-like strength. Our three principal values span the last
of these, and we add a fourth run at the central value of \citet{keane2022}.
At $235$~kg\,m$^{-3}$ the recovered contrast is $1.066$, against $0.9999$ at
$250$: the sense of the answer reverses between the two, inside the published
$1\sigma$ range and just below the bound of \citet{spencer2020}. A linear
extrapolation from the three higher runs would have placed the contrast at
$235$ near $1.006$; the measured value is ten times further from unity, so the
dependence is not linear in this range and the crossing is not predictable
from runs on one side of it. On this body the assumed density, not the data,
sets the direction of the answer (Section~\ref{sec:nulls}).

Four points on the models deserve statement.

First, our Itokawa model encloses $0.0177856$~km$^3$ against the
$(1.84\pm0.092)\times10^{7}$~m$^3$ adopted by \citet{abe2006}, a
difference of $-0.67\sigma$; the bulk density that follows,
$\SI{2012.9}{\kilo\gram\per\cubic\metre}$, is $+0.45\sigma$ from their
$1950\pm140$. The difference is ordinary and within the published
uncertainty, but it is not negligible for our purposes, because
Equation~\eqref{eq:ampsensitivity} amplifies it: a $3.2\%$ offset in
$\densbulk$ moves $\ampfac$ by a factor of order unity, and $\ampfac$ is
the number by which the headline stability ratio must be divided before
any part of it can be called physical. Section~\ref{sec:itokawa:scale}
shows that the contrast and the excess mass move by under one per cent
when the model is rescaled to the published volume, and reports $\ampfac$
for both.

Second, the Arrokoth model of \citet{porter2024} is supplied as two
unjoined shells: $\chi = 4$ over two connected components, with no
repeated or unmatched edges. Geometrically the two lobes touch;
topologically they are separate surfaces. This is why the manifold
criterion of Section~\ref{sec:methods:verify} is stated per component,
and why a plane placed beyond the neck yields a region legitimately in
two pieces.

Two consequences of that arrangement were measured rather than assumed.
The shells touch but barely interpenetrate: a containment test on
$2\times10^{5}$ uniformly distributed samples places $0.044\%$ of the
volume inside both shells. With $88$ samples falling in the intersection
the Poisson standard error is $0.005\%$, so the overlap is resolved
rather than consistent with zero; it is nevertheless small enough that
the double-counted mass is $4\times10^{-4}$ of the total and affects no
reported digit.
% NOTE: an exact boolean mesh intersection would remove the need for this
% qualification and is preferable if the geometry kernel allows it.
Separately, $455$ facets, carrying $0.943\%$ of the surface area, have
their centroids inside the other shell. Those facets are not on the true
surface and have no business in an average taken over it. All Arrokoth
results reported here exclude them, with $\Ubar$ and $\sum_i A_i$
recomputed over the retained set; Section~\ref{sec:nulls:arrokoth} gives
the size of the effect, which is negligible at most planes and decisive
at one.

Third, three shape models of 67P were available. The two that reproduce
their own published volumes to better than $1\%$ agree closely with each
other, while a third, a decimated product, encloses $3.9\%$ less volume
and is not used. The SPC~2017 model is not a closed surface as supplied,
carrying one repeated and two unmatched directed edges; the defect is
inherited by any region cut from it, and is reported where that model is
used.

Fourth, meshes above $50{,}000$ facets were reduced by quadric edge
collapse with the link condition enforced, which preserves the manifold
property and returns the requested facet count exactly. Papers~III
and~IV used vertex clustering, which is equally accurate in the fields it
produces---tested against a homogeneous ellipsoid in
Section~\ref{sec:discussion:ellipsoid}---but does not preserve topology.
The change of method alters the numbers where the reduction is severe: at
a $16\%$ reduction on Itokawa the recovered contrast moves by $0.005\%$,
at $68\%$ on Bennu by $0.054\%$, and at $96\%$ on 67P by $-3.9\%$. The
last figure is larger than the dispersion improvement reported for 67P in
Section~\ref{sec:nulls}, and we carry it into the error budget rather
than record it in passing: on that body the decimation, not the
inversion, dominates.
% Confirmed independent: the -3.875 % is the change in recovered contrast
% between vertex clustering and quadric collapse at 96 per cent reduction on
% the model actually used; the 3.9 % is the volume deficit of a different,
% discarded model. The agreement to two figures is coincidental.

\subsection{Reproducibility}
\label{sec:methods:repro}

The self-test of Section~\ref{sec:methods:verify}, including the
mass-balance audit over all thirteen Itokawa planes, is released as a
runnable script, and the identical values were reproduced on a second
machine: the results reported here were computed under Windows~10
(build~19045) with Python~3.13 and \texttt{polyhedral-gravity}~3.3.1, and
the geometry self-test reproduces its printed values digit for digit under
Ubuntu~24.04 with Python~3.12.3. All code, the driver
scripts and the derived tables are archived at
\href{https://doi.org/10.5281/zenodo.23003040}{doi:10.5281/zenodo.23003040};
the shape
models are not redistributed and are identified by their published
references in Section~\ref{sec:methods:shapes}, with the archive holding
each one recorded in the deposit. The scan intervals of
Section~\ref{sec:methods:minimisation} are recorded per plane in the
released configuration, so that the widening described there can be
audited rather than taken on trust.
% DOI reserved 2026-09-28: 10.5281/zenodo.23003040. Confirm before
% submission that the published record is the one that produced these numbers.